\begin{tikzpicture}
  \tikzset{edge from parent/.style={draw, edge from parent path={(\tikzparentnode.south)-- +(0,-8pt)-| (\tikzchildnode)}}}
  \tikzset{level distance=40pt}
  \tikzset{sibling distance=-4pt}
  \tikzset{every tree node/.style={align=center, anchor=north}}
  \tikzset{every node/.style={font=\footnotesize}}
  \Tree
  [.体系审查
    [.客观
      [.危害行为
        [.没有
          [.沙漠稻草人
            [.对象不能犯 ] % 注意右方括号左边有个空格
          ]
        ]
        [.有 ]
      ]
      [.危害结果 ]
    ]
    [.{主观\\（有无有犯罪故意）}
      [.{没有\\（缺乏犯罪故意型认识错误）}
        [.事实认识错误
          [.{焦点\\过失犯罪 \hz 意外事件} ]
        ]
        [.法律认识错误
          [.违法性的认识可能性 ]
        ]
      ]
      [.{有\\（具有犯罪故意型认识错误）}
        [.预定目标和实害对象法益性质种类
          [.{具体的错误\\（同一犯罪构成内的错误）}
              [.故意
                [.对象错误 ]
              ]
              [.过失
                [.打击错误
                  [.具体符合说 ]
                  [.法定符合说 ]
                ]
              ]
              [.因果关系
                [.狭义的 ]
                [.结果的推迟发生 ]
                [.结果的提前实现 ]
              ]
          ]
          [.{抽象的错误\\（不同犯罪构成间的错误）}
            [.分别分析法 ]
          ]
        ]
      ]
    ]
  ]
\end{tikzpicture}